\section{许华哲的离开：算法创新与价值链壁垒}

机器人大脑之后，访谈进入一个敏感但高价值的问题：联合创始人离开说明了什么？这一章不做八卦判断，而把它当成理解机器人公司壁垒排序的窗口。

访谈录制时，星海图联合创始人许华哲即将离开。主持人追问：这是否意味着现阶段机器人公司算法创新不重要？高继扬的回答很清楚：算法创新重要，但不能脱离整条具身智能价值链独立存在。

\begin{figure}[H]
\centering
\includegraphics[width=0.96\textwidth]{embodied-ai-value-chain.png}
\caption{具身智能价值链与传播周期：整机、客户、数据、AI infra、算法/模型。}
\end{figure}

\begin{knowledgebox}{读图：传播周期决定短期壁垒}
图中传播周期来自访谈中的判断：整机和供应链约 12--18 个月，客户渠道至少 6 个月，数据体系要在整机基础上再加 6--12 个月；而一线团队追赶算法论文或开源方法可能只需 2--3 个月。因此短期壁垒不只看算法，而要看整条链条。
\end{knowledgebox}

\subsection{算法为什么仍重要}

前面用传播周期说明算法不一定是最厚壁垒，但这并不意味着算法可以忽略。更准确的理解是：算法必须嵌入更长周期的资产中，才能把创新变成壁垒。

高继扬并没有贬低算法。他强调星海图算法团队能力很强，也会持续创新。但在当下开源和论文传播速度很快的环境里，算法本身的可复制性更强；如果没有整机、数据、infra、渠道和客户价值支撑，算法创新很难形成长期防守。

\begin{importantbox}{务实创新}
“务实创新”不是不创新，而是先判断创新在整条价值链中的 ROI：它对长期战略有多大贡献，对短期收益有多大贡献，是否能被整机、数据和客户闭环放大。
\end{importantbox}

\subsection{价值观与组织取舍}

算法和壁垒讲完后，问题会自然落到组织：当人、方向和资源不再完全匹配，创始团队如何取舍？这也是“务实创新”真正接受考验的地方。

许华哲离开也被高继扬解释为组织阶段变化：公司在不同阶段需要不同类型的人创造价值，也需要实事求是地做调整。他把价值观定义为两类取舍：面对方向选择时选什么不选什么，面对利益分配时分给谁不分给谁。这比“谁对谁错”更接近创业组织的真实逻辑。

\begin{warningbox}{不要把联合创始人离开简单归因}
创始团队变化可能来自方向、阶段、职责、价值观和个人创业意愿的组合。对外部观察者来说，更重要的是看调整后组织是否继续产出结果，而不是只看人事事件本身。
\end{warningbox}

\subsection{本章小结}

许华哲离开的段落，实际上把具身智能公司的壁垒讲清楚了：算法重要，但算法必须嵌入整机、供应链、数据、infra、模型、分销和客户价值链条。

